%% kapa: exam=1 round=28 year=2017 period=1 form=A subject=civil_law subject_ko=민법 seq=28 number=28
%% source: Q-Net 감정평가사 기출문제 (cst00309) artlSeq=5199603 file=extracted/1008/5199603/2017년 제28회 감정평가사 1차시험 문제지/1교시 A형.pdf
%% conversion: kordoc 4.24.51 -> tools/gen1.py (auto); review=proofread-sonnet; flags=none
\begin{question}{28}
집합건물의 소유 및 관리에 관한 법률상 공용부분에 관한 설명으로 옳지 않은 것은? (다툼이 있으면 판례에 따름)
\begin{choices}
  \item 공용부분은 취득시효에 의한 소유권 취득의 대상이 되지 않는다.
  \item 구조상의 공용부분에 관한 물권의 득실변경은 별도로 등기를 하여야 한다.
  \item 공용부분 관리비에 대한 연체료는 특별한 사정이 없는 한 특별승계인에게 승계되는 공용부분 관리비에 포함되지 않는다.
  \item 관리인 선임 여부와 관계없이 공유자인 구분소유자가 단독으로 공용부분에 대한 보존행위를 할 수 있다.
  \item 어느 부분이 공용부분인지 전유부분인지는 구분소유자들 사이에 다른 합의가 없는 한 그 건물의 구조에 따른 객관적인 용도에 의하여 결정된다.
\end{choices}
\end{question}
\answer{2}
